% =====================================================================================================================
\section{Introduction}
\label{sec:intro}

A researcher finishes a set of experiments, opens a coding agent in the repository, and asks it to write the paper.
The agent reads the result files, drafts the sections, and then revises the draft in a loop with an LLM reviewer until
the score stops rising. Whole-paper systems already work this
way~\cite{lu2024aiscientist,yamada2025aiscientistv2,song2026paperorchestra,meng2026scientistone}, and the reviewer in
the loop is often itself a simulated peer-review system~\cite{jin2024agentreview,goyal2026scholarpeer}. In the third
round of revision, the agent rewrites a results paragraph, and an accuracy of 81.8\% in the data becomes 84.8\% in
the text. No one changed it on purpose. The sentence reads well, the table two pages later still shows the true
value, and the reviewer's next score is no lower.

Readers of the finished paper trust every number in it to follow from the data. A wrong number in a results paragraph
is the error a co-author is least likely to catch and a later paper is most likely to cite. Yet no step in the loop
above compares a typed number with the data behind it.

Existing checks run after the writing is done. Either the reviewer is expected to notice the wrong number, or an
audit agent recomputes the numbers of the finished
paper~\cite{chen2026fabscore,shen2026reagent,autoresearchclaw2026,meng2026scientistone}. Both run long after the
error is made, which happens in a single edit, and both make assumptions about what errors look like. LLM reviewers
raise integrity concerns and accept the paper anyway~\cite{badscientist2026}. An audit that matches printed numbers
against stored values has a different problem. A revising agent computes new analyses from raw files that the
repository never summarized. A lookup then flags this correct work, and it accepts wrong numbers that happen to land
near a stored value. Software engineering solved the same problem decades ago by checking code where it is written:
the compiler catches a type error before a colleague ever reads the pull request.

We ask where a number check should run in a writing agent's loop. To answer, we measure every placement on the same
agents and repositories. To make checks at the edit possible, we built \tool{} (Fig.~\ref{fig:overview}). For every
number it adds or changes, the agent leaves a short recipe. The recipe is a tag in the \LaTeX{} source plus a Python
function that reads the repository's raw data and returns the unrounded value, with a one-line statement of what it
computes. A Claude Code hook reruns the recipe at the edit and tests whether the printed number follows from it at
the printed precision. The same check can run as a tool the agent may call, as feedback
after every edit, or as a gate that refuses the edit. In a generate mode, the build fills in the number from the
recipe, so the agent never types it. The check is deterministic, makes no model call, and on most papers takes under a second.

Measuring placements needs a testbed in which the right answer is known, so we assembled \bench{}. It has four parts.
The first is 60 arXiv papers posted after the models' training cutoff, each in five versions with planted and deleted
numbers. The second is 10 research repositories with public raw results, in which agents write and revise papers
under reviewer pressure. The third is a random sample of 300 numbers those agents wrote, audited blind by two models
and recomputed by hand. The fourth is 24 papers seeded with errors that make the paper's own method look better
(flattering errors), each revised under five placements of the check and two controls.

Reviewers find the planted wrong numbers and quote them in their reviews, but their scores stay where they were:
\emph{the review names the error, but the score does not reflect it}. Writers that revise under review add many
numbers. Almost all of them are correct values computed from raw data. A lookup of stored values and a trace of the
finished paper both miss the few wrong ones, and recomputing the agent's own recipe finds them. Whether the check then
helps depends entirely on where it runs. When it is only offered as a tool to our main writer, as many planted errors
remain as with no check, and how much an offered check removes changes with the writer. When it runs by itself at the
edit, it removes every planted error, for three writer models and for a writer told to keep its results favorable. A
placebo that refuses edits just as often, without checking anything, does no better than no check.

Our core contributions are:
\begin{itemize}
  \item \tool{}: an edit-time recompute check for the numbers a writing agent types. The agent writes a recipe for
        each number, and a hook reruns it against a snapshot of the raw data, with feedback, gate, and generate
        modes. We release it as a Claude Code plugin, a command-line tool, and a git pre-commit hook.
  \item \bench{}: a testbed for number checks in agent-written papers, with post-cutoff papers for reviewers, real
        repositories for writers, a human-recomputed sample of the numbers agents add, and a planted-error protocol
        that separates flattering errors from neutral ones.
  \item A placement study: we show that LLM reviewers of three sizes quote wrong numbers but lower their scores by only
        a few points, that recomputation finds the errors that value lookup and post-hoc tracing miss, and that
        among five placements only those that run at the edit without being asked remove every planted error. Blocking edits adds nothing over
        feedback, even for a writer told to protect its results.
\end{itemize}
